\section{Establish the head width before reading its argument}
The supported heads use the following additional-information values. The major-type
bits remain in the initial byte; they are not part of the argument.

\tagpdfsetup{table/header-rows={1}}
\begin{center}
\begin{tabular}{rrr}
Additional information & Following argument bytes & Complete head bytes \\
\hline
0..23 & 0; value is inline & 1 \\
24 & 1 & 2 \\
25 & 2 & 3 \\
26 & 4 & 5 \\
27 & 8 & 9 \\
\end{tabular}
\end{center}

After reading the initial byte, establish the required head size and compare it with
the unsigned available count before loading any following argument byte. A count of
one permits inspecting \texttt{1b}, but it does not permit reading the eight bytes that
its lower field requests. Reject that truncated head without returning a partial
argument to the caller.

The head reader accepts a supported argument even when its representation is wider
than necessary. Both \texttt{00} and \texttt{18 00} describe unsigned zero. Keeping the
argument separate from its width lets the decoder accept those equivalent inputs,
while the writer later chooses the shortest representation. The value twenty-four in
the lower field is a width selector here, not an inline argument of twenty-four.
